% ===== AGRADECIMENTOS =====
\chapter*{AGRADECIMENTOS}
\thispagestyle{empty}

A construção de um ecossistema com 106 skills, 125 agentes e 41 MCPs que orquestram pipelines científicos de nível Qualis A1 não é obra de um indivíduo isolado, mas sim o produto emergente de uma comunidade distribuída de pesquisadores, desenvolvedores e mantenedores de código aberto.

Agradeço à Secretaria de Educação da Prefeitura Municipal de Crateús pelo apoio institucional que tornou possível esta pesquisa.

Agradecemos à comunidade global de desenvolvedores dos projetos que formam a infraestrutura sobre a qual o OpenCode se ergue: à equipe do \textit{Protocolo MCP} (Model Context Protocol) da Anthropic, que estabeleceu o padrão de interoperabilidade ferramenta-agente; ao ecossistema \textit{Python} e \textit{TypeScript}, que fornecem as linguagens de orquestração; aos mantenedores do \textit{LaTeX} e do \textit{ABNTex2}, que permitem a exportação de documentos com qualidade tipográfica profissional; à comunidade \textit{Qiskit} e \textit{PennyLane} pela instrumentação quântica; e aos incontáveis contribuidores anônimos do \textit{PyPI}, \textit{npm} e \textit{GitHub} cujas bibliotecas compõem a malha de dependências do ecossistema.

Agradecimento especial aos pesquisadores cujos trabalhos fundamentam as disciplinas de engenharia aqui aplicadas: Kent Beck (TDD, 2003), por estabelecer o ciclo RED-GREEN-REFACTOR que estrutura nossa validação; os engenheiros da OpenAI (\textit{Harness Engineering}, 2025), pelo paradigma de \textit{agent harness}; e os autores do \textit{BettaFish} e \textit{MiroFish} (GHJ, 2024-2025), cuja arquitetura OASIS inspirou o pipeline multiagente P14-P18.

Por fim, agradecemos a todos os usuários do ecossistema OpenCode que, com suas demandas e \textit{feedback}, impulsionaram 16 ciclos evolutivos documentados --- cada um elevando o patamar de qualidade do sistema de 85/100 para 99/100.

\vspace{1cm}
\begin{flushright}
  \textit{Crateús, 7 de junho de 2026.}
\end{flushright}
\newpage
